\documentclass[11pt]{article}

\usepackage[margin=1in]{geometry}
\usepackage{amsmath}
\usepackage{amssymb}
\usepackage{booktabs}
\usepackage[colorlinks=true,linkcolor=blue,urlcolor=blue]{hyperref}

\setlength{\parindent}{0pt}
\setlength{\parskip}{0.55em}

\newcommand{\pts}[1]{\hspace{0.35em}\textnormal{\textbf{[#1 points]}}}
\newcommand{\St}{S_T}
\newcommand{\Fzero}{F_0}
\newcommand{\Fone}{F_1}
\newcommand{\Ftwo}{F_2}
\newcommand{\K}{K}
\newcommand{\Var}{\operatorname{Var}}
\newcommand{\Cov}{\operatorname{Cov}}

\title{S181M116 Derivatives and Alternative Investments\\
Assignment 1: Derivative Basics, Futures Markets, Hedging, and Forward Pricing}
\author{Based on Lecture 1 and Lecture 2, Sections 1--3}
\date{}

\begin{document}
\maketitle

\textbf{Instructions.} Submit one PDF file on Moodle. Show your calculations and
briefly explain the economic intuition behind numerical answers. Unless stated
otherwise, ignore transaction costs, taxes, bid-ask spreads, and default risk.
Use continuous compounding whenever an exponential formula is required. Round
monetary answers to two decimals and hedge ratios to four decimals.

\textbf{Total: 20 points.} This assignment covers Lecture 1 and the first three
sections of the Lecture 2 slides: hedge mechanics, basis risk, and forward/futures
pricing. The later Lecture 2 material on commodity behavior, stock-index futures
and beta hedging, and stack-and-roll hedging is not assessed here.

\section*{Question 1: Concepts, Market Structure, and Risk Controls \pts{3}}

\begin{enumerate}
    \item For each item below, state whether it is a derivative. If it is a
    derivative, identify the underlying variable.
    \begin{enumerate}
        \item A contract that pays $\max(\St-75,0)$ in three months.
        \item A share of common stock.
        \item A six-month agreement to exchange EUR 1 million for USD at a fixed
        exchange rate.
        \item A weather contract that pays when average July temperature exceeds
        a fixed threshold.
    \end{enumerate}
    \item A mining company wants to lock in a customized sale price for a
    non-standard grade of metal to be delivered in nine months. Would an
    exchange-traded futures contract or an OTC forward contract be the natural
    starting point? Explain one benefit and one risk of your choice.
    \item In one or two sentences, explain why separating front office, middle
    office, and back office functions matters for derivatives risk management.
\end{enumerate}

\section*{Question 2: Payoff, Profit, and Option Logic \pts{3}}

A European call and a European put have the same strike price, $\K=100$. The call
premium is 7 and the put premium is 5. Ignore interest on the option premia.

\begin{enumerate}
    \item Complete the table below for the buyer of each option.

    \begin{center}
    \begin{tabular}{@{}ccccc@{}}
    \toprule
    $\St$ & Call payoff & Call profit & Put payoff & Put profit \\
    \midrule
    88  & & & & \\
    100 & & & & \\
    116 & & & & \\
    \bottomrule
    \end{tabular}
    \end{center}

    \item Explain briefly why an option can have a positive payoff and still
    produce a negative profit.
    \item Compare the buyer's downside in a long call with the downside in a long
    forward on the same underlying.
\end{enumerate}

\section*{Question 3: Futures Settlement, Margin, and Liquidity Risk \pts{3}}

A trader is \textbf{short} six wheat futures contracts. Each contract is for
5,000 bushels. Yesterday's settlement price was 580.25 cents per bushel and
today's settlement price is 592.75 cents per bushel. The margin account balance
before today's settlement was USD 21,000. The maintenance margin is USD 16,000
and the initial margin is USD 21,000.

\begin{enumerate}
    \item Compute today's futures profit or loss.
    \item Compute the margin account balance after daily settlement.
    \item Is there a margin call? If so, how much must be deposited?
    \item Briefly explain why daily settlement can create liquidity risk even
    for a hedge that reduces the firm's underlying price risk.
\end{enumerate}

\section*{Question 4: Hedge Direction, Basis Risk, and Effective Prices \pts{4}}

A firm will sell a commodity in three months. It hedges today by shorting futures
at $\Fone=47.80$. Three months later, it sells the physical commodity in the spot
market at $S_2=45.90$ and closes the futures position at $\Ftwo=46.20$.

\begin{enumerate}
    \item Why is a short hedge the appropriate hedge direction for this exposure?
    \item Compute the gain or loss on the futures position per unit.
    \item Compute the effective sale price per unit.
    \item Compute the basis at hedge close, $b_2=S_2-\Ftwo$, and verify that the
    effective sale price equals $\Fone+b_2$.
    \item Suppose the firm originally expected the closing basis to be $-0.60$.
    Was the realized effective sale price better or worse than expected? Explain.
\end{enumerate}

\section*{Question 5: Cross Hedging and the Minimum-Variance Hedge Ratio \pts{4}}

An airline will buy 1,680,000 gallons of jet fuel in three months. It cross
hedges with heating-oil futures because the futures contract is more liquid.
Each futures contract covers 42,000 gallons. Historical data give
\[
    \rho=0.82,\qquad \sigma_S=0.24,\qquad \sigma_F=0.20,
\]
where $\Delta S$ is the change in the jet-fuel spot price and $\Delta F$ is the
change in the heating-oil futures price.

\begin{enumerate}
    \item Compute the minimum-variance hedge ratio
    \[
        h^*=\rho\frac{\sigma_S}{\sigma_F}.
    \]
    \item Compute the number of futures contracts implied by the hedge ratio.
    Should the airline round up or down? Give a practical recommendation.
    \item Should the airline go long or short futures? Explain.
    \item Name two sources of residual risk that remain after this hedge.
\end{enumerate}

\section*{Question 6: Forward Pricing and Contract Value \pts{3}}

\begin{enumerate}
    \item A stock has current price $S_0=96$. The present value of known
    dividends paid during the life of a nine-month forward contract is 1.80. The
    continuously compounded risk-free rate is 4\% per year. Compute the
    no-arbitrage forward price.
    \item A dealer quotes a forward price of 99.50 for the same contract. Is the
    quote too high or too low relative to no-arbitrage? Describe the arbitrage
    trade and compute the locked-in profit per share.
    \item You are long an old six-month forward contract with delivery price
    $\K=92$. A new forward contract with the same maturity has delivery price
    $\Fzero=94$. The continuously compounded risk-free rate is 4\% per year.
    Compute the current value of the old long forward.
\end{enumerate}

\end{document}
